\subsection{Rank one and the distinction between finite and geometric tests}
\begin{statement}{A.44}{Finite kernel refinements of an elliptic section}{Proved}
Let $E/\Q$ be an elliptic curve, $P\in E(\Q)$ a point of infinite order,
$D\subset E$ a finite reduced divisor, and $U=E-D$. Choose the origin of
$E$ and a geometric base point, so that
$\pi_1(E)=T E\rtimes G_\Q$, where $T E=\prod_\ell T_\ell E$.
For $a\in\widehat\Z$, write $s_a$ for the section with Kummer class
$a\kappa(P)$. Suppose that a continuous quotient of arithmetic groups
$\pi_1(U)\twoheadrightarrow\Gamma\twoheadrightarrow\pi_1(E)$
has finite kernel in its second arrow. Then every $s_a$ lifts to a
continuous section of $\Gamma\to G_\Q$. The same holds for any finite
collection of these refinements simultaneously.
\end{statement}
\begin{proof}
Put $K=\ker(\Gamma\to\pi_1(E))$. Choose an open normal subgroup
$J\triangleleft\Gamma$ with $J\cap K=1$, and put $V=\operatorname{im}J$.
The latter is open and normal, and the natural continuous homomorphism
\[
 \Gamma\longrightarrow
 \pi_1(E)\mathop{\times}_{\pi_1(E)/V}\Gamma/J
\]
is an isomorphism. Injectivity is exactly $J\cap K=1$. For surjectivity,
lift a compatible pair first from $\Gamma/J$ to $\Gamma$; its error in
$\pi_1(E)$ belongs to $V$ and is corrected by an element of $J$.
The map is then a homeomorphism by compactness.

Choose an integer $M>0$ annihilating the image of $T E$ in
$\pi_1(E)/V$. The finite representation induced by $s_a$ depends only
on $a$ modulo $M$. Choose an integer $m\equiv a\pmod M$ with $mP\notin D$;
this is possible because $P$ has infinite order and $D$ is finite.
The rational point $mP$ gives a section of $\pi_1(U)$ and hence of
$\Gamma/J$. After choosing its path so that its compact Kummer cocycle
is $m\kappa(P)$, this finite section and $s_a$ agree in $\pi_1(E)/V$.
Their pair under the displayed isomorphism is the required lift.
For several refinements, use their common quotient of $\pi_1(U)$;
its kernel over $\pi_1(E)$ embeds in the finite product of their kernels.
\end{proof}

For example, an element $a\in\widehat\Z$ whose $2$-adic component is
$0$ and whose odd components are $1$ is not an integer, and passes
every test in A.44. This statement does not assert that these lifts
are Selmer, or compatible as the finite refinements vary. Nor does
it retain the full fundamental group of a punctured congruence cover.
The comparison is necessary: a finite group such as $A_5$ and the
fundamental group of its corresponding punctured cover are different
objects, and have different lifting problems.

\begin{statement}{A.45}{The remaining rank-one lifting problem}{Open}
For the least-conductor cubic twist of analytic rank one in A.43,
the complete sets requested in task A.12 remain undetermined: the
classes $a\in\widehat\Z$ lifting to the intrinsic full profinite
metabelian quotient, and to the geometric quotients retaining the
punctured covers over $X(3),X(4),X(5),X(7)$, with the Selmer conditions
required by Lead 5. No assertion that all $a\notin\Z$ die in this
finite list of geometric tests is established here.
\end{statement}
The rank-zero argument A.24 cannot be applied to the rank-one
Jacobian; its finite-descent set already contains the closure of
$E(\Q)$. The finite-kernel approach gives A.44 and hence cannot
separate nonintegral exponents. The next unresolved step is an
explicit calculation of the nonabelian obstruction above this
closure in one of the stated geometric quotients, together with
its local conditions and its continuity condition. A computation
of the rational points of selected twists of $X(7)$ alone does not
supply this comparison for every profinite exponent.
